% introduction.tex — HapGraph Introduction (v2 — post Phase B review)
% Fixes: GLOBETROTTER citation, AdmixtureBayes timing clarification,
%        "well-calibrated" -> specific numbers, T-not-estimated for 1kGP,
%        S2 over-detection caveat, "incompatible" -> "separate"

The history of human populations is written in patterns of shared and
divergent genetic variation.
When two populations exchange migrants — whether through conquest, trade, or
geographic contact — their descendants carry a mosaic genome encoding
both the identity of the contributing ancestral groups and the timing of
the exchange.
Reconstructing this history requires inferring not only
\emph{who mixed with whom} (the admixture graph topology and proportions)
but also \emph{when} (the admixture timing in generations before the present).
These two questions are today addressed by separate frameworks applied
sequentially, with no principled joint inference; this gap is the
motivation for the present work.

\paragraph{Admixture graph inference.}
The admixture graph — a directed acyclic graph in which internal nodes
represent ancestral populations and admixture nodes represent lineage mixtures
— is the standard model for joint representation of population divergence and
admixture \citep{Reich2009,Patterson2012}.
\citet{Patterson2012} established that F-statistics provide a rich,
statistically tractable summary of admixture graph structure:
negative $F_3(C;\,A,B)$ identifies admixed populations, and $F_4$ statistics
constrain topology and proportions.
\treemix{} \citep{Pickrell2012} operationalised this into the first automated
graph inference tool, using greedy maximum-likelihood search on the $F_2$
covariance matrix.
\textsc{ADMIXTOOLS2} \citep{Maier2023} and \textsc{AdmixtureBayes}
\citep{Nielsen2023} extended the framework with formal Bayesian topology
search and posterior distributions over graph parameters.
These tools infer graph topology and admixture proportions with rigour,
but they do not jointly model the temporal dimension of admixture.

\paragraph{Admixture timing.}
Haplotype-based methods exploit the fact that admixture introduces ancestral
segments whose lengths decay over time as recombination fragments them.
\alder{} \citep{Loh2013} models the decay of weighted linkage disequilibrium
to estimate $T$ for events $\gtrsim 5$ generations ago.
\textsc{GLOBETROTTER} \citep{Hellenthal2014} takes a complementary approach,
inferring both admixture sources and timing from LD patterns; however, it
operates on a source--target model rather than a full graph topology and does
not incorporate identity-by-descent (IBD) statistics.
More recently, IBD-based approaches have exploited the observation that the
expected mean length of inter-population IBD segments decays as $100/T$~cM,
providing a direct genomic clock on admixture timing
\citep{Browning2020,Palamara2012}.
Despite this complementarity, current practice couples the two estimation
problems only sequentially: users infer topology with F-statistic tools and
subsequently run a timing method under the fixed topology, propagating
topology uncertainty into timing estimates without any feedback loop.

\paragraph{The case for joint inference.}
Topology misspecification — a common outcome of greedy graph search — biases
downstream timing estimates, because the IBD likelihood depends on which
populations are designated as admixed and which as sources.
Conversely, timing information can resolve topology degeneracies: two competing
graphs may be equally consistent with $F_2$ statistics yet make different IBD
predictions, allowing the IBD likelihood to discriminate between them.
A jointly calibrated framework would propagate uncertainty coherently across
topology, proportions, and timing, yielding credible intervals rather than
the point estimates returned by current pipelines.

\paragraph{This work: \hapgraph{}.}
We introduce \hapgraph{}, a Python package for joint admixture graph inference
using F-statistics for topology and proportion estimation (Stage I) followed
by Bayesian IBD-based timing inference (Stage II).
The two stages are deliberately decoupled: preliminary experiments showed
that incorporating IBD into the topology search causes systematic edge
over-detection; IBD is therefore reserved exclusively for timing inference
after the topology is fixed.
Applied to coalescent simulations under three scenarios (S1: pure tree,
$K=0$; S2: one cross-clade event, $K=1$; S3: two independent events, $K=2$),
\hapgraph{} achieves a false positive rate of 0\% on pure trees and
recovers all true admixture target populations with 100\% recall.
Parameter estimation under oracle topology achieves alpha MAE of
$0.054$--$0.061$ and T MAE of $10.9$--$11.4$ generations; the 95\% credible
intervals for $T$ achieve 85--93\% empirical coverage meeting the pre-specified
target, while alpha CIs remain conservative at 55--62\% (see Discussion).
Applied to the 1000 Genomes Project 2022 high-coverage dataset
\citep{ByrskaBishop2022}, \hapgraph{} infers $K=8$ admixture events whose
targets and proportions are consistent with published population history
across African American, Latino, and South Asian populations
\citep{Bryc2015,Moreno-Estrada2013,Moorjani2013};
admixture timing ($T$) was not estimated for the real data application,
as it requires pre-computation of IBD segments by \hapibd{} and is
deferred to future work.

\paragraph{Contributions.}
(i) A two-stage pipeline integrating F-statistics and IBD sharing for joint
admixture graph topology inference and timing estimation;
(ii) a topology-independent $F_3$ method-of-moments estimator for admixture
proportions, robust to neighbor-joining misplacement;
(iii) a BIC-based greedy search with a scale-adaptive stopping criterion
achieving 0\% false positive rate on pure trees and 100\% target recall
across $K \leq 2$ simulation scenarios;
(iv) an open-source Python implementation with documented validation on
coalescent simulations and application to the 1kGP 2022 dataset.
